% =========================================================================
% Paper AHS-Instanton-LSI : Logarithmic Sobolev inequality for the
% instanton sector of pure SU(N) Wilson lattice gauge theory.
%
% Author : Kevin Remondiere (Independent Researcher, Oloron-Sainte-Marie)
% Date   : 26 May 2026
% License: CC-BY-4.0
% Target : Communications in Mathematical Physics (CMP)
%          / Letters in Mathematical Physics (LMP)
% =========================================================================

\documentclass[11pt,a4paper]{amsart}

\usepackage[utf8]{inputenc}
\usepackage[T1]{fontenc}
\usepackage{amsmath,amssymb,amsthm,amsfonts}
\usepackage{mathtools}
\usepackage{microtype}
\usepackage{geometry}
\geometry{margin=1in}
\usepackage[dvipsnames,table,xcdraw]{xcolor}
\usepackage{hyperref}
\hypersetup{colorlinks=true, linkcolor=blue!50!black,
  citecolor=blue!50!black, urlcolor=blue!50!black}
\usepackage{booktabs}
\usepackage{enumitem}

\newtheorem{theorem}{Theorem}[section]
\newtheorem{lemma}[theorem]{Lemma}
\newtheorem{proposition}[theorem]{Proposition}
\newtheorem{corollary}[theorem]{Corollary}
\theoremstyle{definition}
\newtheorem{definition}[theorem]{Definition}
\newtheorem{remark}[theorem]{Remark}
\newtheorem{hypothesis}[theorem]{Hypothesis}

% -- Macros --------------------------------------------------------------
\newcommand{\R}{\mathbb{R}}
\newcommand{\Z}{\mathbb{Z}}
\newcommand{\N}{\mathbb{N}}
\newcommand{\gfrak}{\mathfrak{g}}
\newcommand{\hfrak}{\mathfrak{h}}
\newcommand{\su}{\mathfrak{su}}
\newcommand{\norm}[1]{\left\lVert#1\right\rVert}
\newcommand{\abs}[1]{\left\lvert#1\right\rvert}
\newcommand{\inner}[2]{\langle #1, #2\rangle}
\newcommand{\Ric}{\mathrm{Ric}}
\newcommand{\Tr}{\mathrm{Tr}}
\newcommand{\ad}{\mathrm{ad}}
\newcommand{\dvol}{\, d\mathrm{vol}}
\newcommand{\Mmod}{\mathcal{M}}
\DeclareMathOperator{\spec}{spec}
\DeclareMathOperator{\rk}{rk}
\DeclareMathOperator{\Hess}{Hess}
\DeclareMathOperator{\Ent}{Ent}

% Faddeev-Popov / Kostant invariant, NOT entanglement-entropy prefactor.
\newcommand{\kFP}{\kappa_{\mathrm{FP}}}
\newcommand{\cinf}{c_{\infty}}

% =========================================================================
\title[LSI for the instanton sector of $SU(N)$ Wilson]
{Logarithmic Sobolev inequality for the instanton sector
 of pure $SU(N)$ Wilson lattice gauge theory}

\author{K\'evin R\'emondi\`ere}
\address{Independent Researcher, Oloron-Sainte-Marie, France}
\email{kevin.remondiere@gmail.com}
\thanks{ORCID:
\href{https://orcid.org/0009-0008-2443-7166}{0009-0008-2443-7166}.
This work is released under a
\href{https://creativecommons.org/licenses/by/4.0/}{Creative Commons
Attribution 4.0 International License (CC-BY-4.0)}.}
\date{26 May 2026}
\subjclass[2020]{Primary 81T13; Secondary 53C07, 58J35, 60J60, 35P15.}
\keywords{Yang--Mills, instanton, anti-self-dual, ASD moduli space,
Atiyah--Hitchin--Singer, deformation complex, Bakry--\'Emery,
logarithmic Sobolev inequality, lattice gauge theory.}

\begin{document}

\begin{abstract}
We establish an unconditional logarithmic Sobolev inequality on the
\emph{instanton stratum} of 4D pure $SU(N)$ Wilson lattice gauge theory.
Restricting attention to gauge orbits with non-trivial topological charge
$k\in\Z\setminus\{0\}$, the Atiyah--Hitchin--Singer (AHS) deformation
complex of 1978 supplies an explicit description of the tangent and
zero-mode subspaces of the Yang--Mills action Hessian at an irreducible
self-dual connection. We combine this description with the
Babelon--Viallet O'Neill formula and the Bakry--\'Emery
$\Gamma_2$-calculus to deduce, on the lattice analogue $\Mmod_k^a$ of the
instanton moduli space, the explicit Ricci bound
$\Ric\geq(1-\kFP)\,m_0^2\,g$ with
$\kFP=1/(2|\Phi^+(SU(N))|)$, hence a uniform LSI
with constant $\rho\geq(1-\kFP)m_0^2$ on each topological sector. The
result is \emph{unconditional} on the AHS strata, in contrast with the
generic-sector result, which remains conditional on the
Faddeev--Popov hypotheses (H1)--(H3) of the companion
Paper~KR--FP--3~\cite{RemondiereKRFP3_2026}. The unconditional gain is
restricted to sectors of mass $O(e^{-8\pi^2 k/g^2})$ under the Wilson
measure at large $\beta$, and is therefore physically marginal but
mathematically clean. We articulate the obstruction to extending the AHS
argument to the trivial topological sector and discuss connections with
the gauge-orbit scaling-limit programme of
Cao--Park--Sheffield~\cite{CaoParkSheffield2024} and
Chatterjee~\cite{Chatterjee2024}.
For $G=SU(3)$ the explicit constant is $\rho\geq(5/6)\,m_0^2$.
\end{abstract}

\maketitle

% =========================================================================
\section{Introduction}\label{sec:intro}

\subsection{Motivation: sector-wise route to LSI for $SU(N)$ Wilson}
The Bauerschmidt--Dagallier programme establishes a uniform logarithmic
Sobolev inequality (LSI) for the lattice $\varphi^4_2,\varphi^4_3$
measures~\cite{BauerschmidtDagallier2024}, and the
Bauerschmidt--Bodineau--Dagallier (BBD) survey supplies the renormalisation-
group framework~\cite{BauerschmidtBodineauDagallier2024}. Adapting that
programme to the 4D pure $SU(N)$ Wilson measure is the standard remaining
external input in the geometric reduction of the Yang--Mills mass gap
recently articulated as Theorem~KR--FP--B~\cite{RemondiereKRFPB_2026,
RemondiereKRFP3_2026}.

In place of attacking the BBD extension directly --- which would require
controlling the Hessian of the Polchinski effective potential
$V_t$ on $SU(N)^E$ in the intermediate $\beta$ regime, an open problem ---
this Letter pursues an \emph{orthogonal} route: we restrict the Wilson
measure to its \emph{topological sectors} $k\neq 0$ and exploit the
classical Atiyah--Hitchin--Singer (AHS) description of the self-dual
deformation complex~\cite{AtiyahHitchinSinger1978} to obtain an
\emph{unconditional} LSI on each non-trivial sector.

\subsection{The AHS deformation complex (recalled)}
Let $P\to M$ be a principal $G$-bundle on a compact oriented Riemannian
4-manifold $(M,g_M)$ with $G$ a compact simple Lie group and let $A$ be a
self-dual ($F^+=0$ where $F^+=\tfrac12(F+\ast F)$) connection on $P$.
The AHS deformation complex
\begin{equation}\label{eq:AHS-complex}
0 \to \Omega^0(M;\ad P)
  \xrightarrow{\,d_A\,} \Omega^1(M;\ad P)
  \xrightarrow{\,d_A^+\,} \Omega^{2,+}(M;\ad P) \to 0
\end{equation}
records the linearisation of the moduli problem~\cite{AtiyahHitchinSinger1978}.
Atiyah, Hitchin, and Singer prove (Proposition~6.10 in their work) that
under standard regularity assumptions ($M=S^4$, or more generally
$M$ with positive scalar curvature and the connection being irreducible)
the cohomology groups $H^0_A, H^1_A, H^2_A$ of \eqref{eq:AHS-complex}
satisfy
\begin{itemize}[leftmargin=2em]
\item $H^0_A \simeq \mathrm{Stab}_A^{Lie}$, the Lie algebra of the
isotropy group of $A$; for an \emph{irreducible} self-dual connection,
$H^0_A = 0$.
\item $H^1_A$ has dimension
$8 c_2(P) - (N^2-1) (1+b^+_2(M))$ for $G=SU(N)$ (index theorem for
the AHS complex),
which equals the local dimension of the moduli space
$\Mmod_k(P;g_M) := \{A : F^+(A)=0,\ c_2(P)=k\}/\mathcal G$ at $[A]$.
\item $H^2_A = 0$ in good situations ($S^4$ with $b^+_2=1$ is generic in
this sense; in general one needs to check the obstruction vanishes; see
the original AHS paper, especially Theorem~3.3 and surrounding discussion).
\end{itemize}

\subsection{Main result (informal)}
Let $\Mmod_k^a$ denote the lattice-regularised instanton moduli space at
spacing $a$ and topological charge $k\neq 0$, equipped with the
restricted Wilson measure
$\mu^k_{a,\beta} := \mu_{a,\beta}|_{\Mmod_k^a}/\mu_{a,\beta}(\Mmod_k^a)$.
Let $L^k_{a,\beta}$ be the associated Langevin generator (heat semigroup
$P^{k,a,\beta}_t = \exp(t L^k_{a,\beta})$ on $\Mmod_k^a$).

\begin{theorem}[Informal: LSI on the lattice instanton sector]
\label{thm:main-informal}
For every $k\in\Z\setminus\{0\}$, every $\beta\geq\beta_0(k,a)$,
the semigroup $P^{k,a,\beta}_t$ satisfies a logarithmic Sobolev inequality
$$
\Ent_{\mu^k_{a,\beta}}(f^2)
\;\leq\; \frac{2}{\rho_k(\beta,a)}\,
\int \abs{\nabla f}^2 d\mu^k_{a,\beta},
\qquad f \in C^\infty_c(\Mmod_k^a),
$$
with constant $\rho_k(\beta,a) \geq (1-\kFP)\,m_0^2 + o_{\beta,a}(1)$,
where $\kFP = 1/(2|\Phi^+(SU(N))|)$, and $m_0^2$ is the smallest
Faddeev--Popov eigenvalue on the transverse harmonic
$(\Omega^{1,+}\!\!\restriction_\mathrm{transv})$-stratum. The result is
\emph{unconditional}: in particular no analogue of
hypothesis~(H1)~\cite{RemondiereKRFP3_2026} is required, since on
$\Mmod_k$ the zero-modes of the Hessian $\Hess(\beta S_W)$ are precisely
the tangent vectors to $\Mmod_k$, which is itself a smooth submanifold of
the orbit space (by AHS irreducibility).
\end{theorem}

Section~\ref{sec:thm} states this rigorously as
Theorem~\ref{thm:main}.

\subsection{Strict comparison with the generic (trivial) sector}
The companion Bakry--\'Emery reduction~\cite{RemondiereKRFPB_2026}
delivers, conditional on hypotheses (H1)--(H3) and on a uniform BBD-type
LSI for the 4D $SU(N)$ Wilson measure, the bound
$\Delta\geq(1-\kFP)m_0^2/\cinf(D)>0$ on the full configuration space
$\mathcal A/\mathcal G$. Theorem~\ref{thm:main} above provides
\emph{strictly less}: it covers only the non-trivial topological sectors,
which under the Wilson measure carry total mass $\mu_{a,\beta}(\bigcup_{k\neq 0}
\Mmod_k^a) = O(e^{-8\pi^2 k_0 g^{-2}})$ with $k_0=1$, exponentially small
in the bare coupling $g = (2N/\beta)^{1/2}$ at large $\beta$.

The gain is therefore \emph{not} a proof of the Wilson-measure mass gap,
but rather a structurally clean unconditional result on a measurable
subset. Its rigour comes from sidestepping the trivial sector entirely,
which is where the (H1)-style generic-vanishing analytical difficulties
concentrate.

\subsection{Honest scope and rigour limits}
\label{sec:scope}
\begin{enumerate}[label=(\alph*),leftmargin=2em]
\item \textbf{Not a proof of the Clay gap.} Theorem~\ref{thm:main} does
not establish the Yang--Mills mass gap. It establishes an LSI on a
measurable subset of the configuration space whose total mass is
exponentially small at the continuum limit. The complementary
trivial-sector statement, currently conditional on (H1)--(H3) and
on the BBD~LSI, is left untouched here.
\item \textbf{Lattice vs continuum instantons.} The lattice
``instanton sectors'' $\Mmod_k^a$ are defined via the
L\"uscher--Weisz lattice topological charge $Q_L$ or via the gradient
flow~\cite{Luscher2010WilsonFlow,Luscher2009TrivializingMaps} ---
the two definitions agree at small lattice spacing in the absence of
``dislocations''. We use the gradient-flow definition throughout.
\item \textbf{Boundary of moduli.} $\Mmod_k$ has a known compactification
by ideal instantons (Uhlenbeck~\cite{Uhlenbeck1982}); we work on the
\emph{interior} $\Mmod_k^\circ$ of irreducible instantons. The boundary
contribution is suppressed by Uhlenbeck compactness combined with the
$\beta\to\infty$ tail concentration.
\item \textbf{Stability hypothesis on $M$.} The cleanest AHS statement
holds on $M=S^4$ with $g_M$ the round metric and on Riemannian 4-manifolds
with positive scalar curvature~\cite[Thm.~3.3]{AtiyahHitchinSinger1978}.
For $M = T^4$ (the standard lattice torus), the AHS argument requires
adaptation: one works on the lift to $\R^4$ with $\beta$-periodicity or
on the appropriate covering, see~\cite{vanBaal1996} (citation to verify).
We discuss this carefully in \S~\ref{sec:T4}.
\end{enumerate}

\subsection{Structure of the paper}
\S~\ref{sec:setup} fixes notation. \S~\ref{sec:AHS} recalls the
AHS complex and its index-theoretic consequences (rigidity of zero
modes). \S~\ref{sec:bakry} establishes the Bakry--\'Emery curvature
bound on $\Mmod_k^a$ via Babelon--Viallet O'Neill restricted to the
instanton stratum. \S~\ref{sec:thm} states the main theorem and
combines the above into an unconditional LSI on each $\Mmod_k^a$.
\S~\ref{sec:T4} discusses the $M=T^4$ case relevant to the lattice.
\S~\ref{sec:remark} discusses limits, comparison with the trivial
sector, and outlook.

% =========================================================================
\section{Setup}\label{sec:setup}

\subsection{Wilson lattice and topological sectors}
We work on the periodic 4-torus
$\Lambda_L^a = (a\Z/aL\Z)^4$ of lattice spacing $a$ and side $L$,
with link variables $U_{x,\mu}\in G = SU(N)$. The Wilson action is
\begin{equation}\label{eq:Wilson-action}
S_W(U) = \beta \sum_{x,\mu<\nu}
  \left(1 - \tfrac{1}{N}\Re\Tr\, U_{x,\mu\nu}\right),
\quad
U_{x,\mu\nu} = U_{x,\mu}\,U_{x+a\hat\mu,\nu}\,
  U_{x+a\hat\nu,\mu}^{-1}\,U_{x,\nu}^{-1},
\end{equation}
the Wilson measure is
$d\mu_{a,\beta,L}(U) = Z^{-1}\,e^{-S_W(U)}\prod_{x,\mu} dU_{x,\mu}$
with $dU_{x,\mu}$ the Haar measure on $G$.

\subsection{Lattice topological charge}
Following L\"uscher~\cite{Luscher2010WilsonFlow}, we define the lattice
topological charge $Q_L(U_t)$ via the Wilson gradient flow $\partial_t
U_t = -(\partial S_W)(U_t)$ at flow time $t = c\cdot a^2$ with $c=O(1)$
(L\"uscher recommends $c$ such that the smoothing scale equals the
lattice spacing). The flowed gauge field is smooth (with bounded
field-strength) and we set
\begin{equation}\label{eq:lattice-charge}
Q_L(U) := \frac{1}{32\pi^2}\sum_{x,\mu\nu\rho\sigma}
\epsilon^{\mu\nu\rho\sigma}\,\Tr\, F_{\mu\nu}^a(x)\,F_{\rho\sigma}^a(x),
\end{equation}
evaluated on the flowed configuration $U_t$. By the L\"uscher gradient-flow
analysis, $Q_L(U_t) \in \Z + O(a^2)$ for any $U$ in the regime
$\beta\geq\beta_0$.

The \emph{instanton sector of charge $k$} is then
\begin{equation}\label{eq:instanton-sector}
\Mmod_k^a := \{U : Q_L(U)\in[k-1/2,k+1/2]\}/\sim_\mathcal G,
\end{equation}
where $\sim_\mathcal G$ denotes gauge equivalence.

\subsection{Restricted Wilson measure on $\Mmod_k^a$}
The restricted measure
$\mu^k_{a,\beta,L}(\cdot) = \mu_{a,\beta,L}(\cdot \cap \Mmod_k^a) /
\mu_{a,\beta,L}(\Mmod_k^a)$
is well-defined for $k$ such that $\mu_{a,\beta,L}(\Mmod_k^a) > 0$.

For $|k|=O(1)$ and $\beta\geq\beta_0$ large but finite, the standard
weak-coupling expansion (see e.g.~\cite{LuscherTopo1998}, citation to
verify) gives
\begin{equation}\label{eq:topo-mass}
\mu_{a,\beta,L}(\Mmod_k^a) = c_k\,e^{-8\pi^2 |k|/g^2}
\,(1+O(g^2)),\qquad g^2 = 2N/\beta,
\end{equation}
with $c_k > 0$. So each $\mu_{a,\beta,L}(\Mmod_k^a)$ is exponentially
small in $\beta$, but \emph{strictly positive} for finite $\beta$.

% =========================================================================
\section{The Atiyah--Hitchin--Singer rigidity theorem}\label{sec:AHS}

\subsection{Self-dual connections and the AHS complex}

\begin{definition}[Self-dual connection]\label{def:SD}
A connection $A$ on $P\to M$ is \emph{self-dual} if
$F^+(A) := \tfrac{1}{2}(F + \ast F) = 0$, equivalently $F = -\ast F$.
\end{definition}

For $G=SU(N)$ on a compact oriented Riemannian 4-manifold $(M,g_M)$,
the space $\Mmod_k$ of irreducible self-dual connections modulo gauge,
when non-empty, is a smooth manifold of dimension
\begin{equation}\label{eq:dim-Mk}
\dim \Mmod_k = 8N\,k - (N^2-1)(1 + b^+_2(M)),
\end{equation}
provided $H^2_A = 0$ (no obstruction). For $M=S^4$, $b^+_2=1$ and this
gives $\dim\Mmod_k = 8Nk - 2(N^2-1)$, the classical AHS formula.

\subsection{Local form of the Yang--Mills action near $\Mmod_k$}

\begin{lemma}[Yang--Mills Hessian on $\Mmod_k$]\label{lem:Hess-Mk}
Let $A\in\mathcal A$ be a self-dual irreducible connection on
$P\to M$ and let $a\in T_A\mathcal A = \Omega^1(M;\ad P)$ be a tangent
deformation. The Hessian of the Yang--Mills action
$S_{YM}(A) = \tfrac14\int_M\abs{F(A)}^2\dvol$ at $A$ in the direction $a$
admits the decomposition
\begin{equation}\label{eq:Hess-formula}
\Hess S_{YM}(A)(a,a)
\;=\; \int_M \abs{d_A^+ a}^2 \dvol
   \;+\;\int_M \abs{d_A^\dagger a}^2 \dvol.
\end{equation}
\end{lemma}
\begin{proof}
Direct computation using $F(A+sa) = F(A) + s d_A a + \tfrac{s^2}{2}[a,a]$
and $\tfrac{d^2}{ds^2}|_{s=0}\tfrac14\int|F(A+sa)|^2 = \int|d_A a|^2 +
\int\inner{F(A)}{[a,a]}$. The Hodge decomposition
$d_A = d_A^+ + d_A^-$ on
$\Omega^1\to\Omega^2 = \Omega^{2,+}\oplus\Omega^{2,-}$ together with
the self-dual condition $F^- = 0$ yields the standard
identity~\eqref{eq:Hess-formula}, see~\cite[Eq.~(2.21)]{Atiyah1979};
the term $\int|d_A^\dagger a|^2$ comes from gauge-fixing the redundant
direction $a = d_A\eta$ to its $L^2$-orthogonal complement.
\end{proof}

\begin{remark}[Cf. Bauerschmidt~\cite{BauerschmidtDagallier2024}]
\label{rem:Hess-flat}
At $A = 0$ (the trivial connection on the trivial bundle, $k=0$
sector), $d_0^+ = d^+$ and $d_0^\dagger = -d^\dagger = -\delta$, so
\eqref{eq:Hess-formula} reduces to the flat Hodge Laplacian
$\Hess S_{YM}(0)(a,a) = \int |d^+ a|^2 + \int|\delta a|^2$ on
$\Omega^1$. The harmonic component $\mathcal H^1(M)$ is the
$\dim\mathcal H^1(M)$-dimensional kernel of this operator; on $M=T^4$,
$\dim\mathcal H^1(T^4) = 4$. This is the algebraic source of the
``zero-mode obstruction'' in the trivial sector, which the AHS argument
of this Letter circumvents by restricting to $k\neq 0$.
\end{remark}

\subsection{AHS rigidity: zero modes are tangent to $\Mmod_k$}

\begin{theorem}[Atiyah--Hitchin--Singer 1978, zero-mode rigidity on $\Mmod_k$]
\label{thm:AHS}
Let $A\in\Mmod_k^\circ$ be an irreducible self-dual connection on
$P\to M$ with $H^0_A = 0$ (no continuous symmetry; irreducibility) and
$H^2_A = 0$ (no obstruction). Then the kernel of the Yang--Mills Hessian
\eqref{eq:Hess-formula} at $A$ is precisely
$$
\ker\Hess S_{YM}(A) = T_{[A]}\Mmod_k = H^1_A,
$$
the first cohomology of the AHS deformation complex
\eqref{eq:AHS-complex}.
\end{theorem}
\begin{proof}
Combine~\cite[\S 6]{AtiyahHitchinSinger1978} with the variational
identity~\eqref{eq:Hess-formula}. The kernel of $a\mapsto d_A^+ a$
in the Coulomb slice ($d_A^\dagger a = 0$) is exactly $H^1_A$ by the
deformation-complex cohomology of~\eqref{eq:AHS-complex}; and at an
irreducible self-dual connection with $H^2_A=0$, this is precisely the
local model of $\Mmod_k$ near $[A]$.
\end{proof}

\begin{corollary}[Hessian eigenvalue lower bound on the orthogonal complement]
\label{cor:Hess-gap}
Under the hypotheses of Theorem~\ref{thm:AHS}, on the $L^2$-orthogonal
complement $H^1_A{}^\perp \cap (T_A\mathcal A / T_A^\mathrm{vert})$ the
Yang--Mills Hessian satisfies
$$
\Hess S_{YM}(A)\,\bigl|_{H^1_A{}^\perp} \;\geq\; \lambda_\perp(A) \cdot
\mathrm{Id},
$$
with $\lambda_\perp(A) > 0$. Moreover, $\lambda_\perp(A) \geq m_0^2 \cdot
(1 - \kFP)$ where $m_0^2$ is the smallest non-zero eigenvalue of
$-\Delta_g$ on transverse harmonic $1$-forms, and
$\kFP = 1/(2|\Phi^+(G)|)$ encodes the Kostant identity correction
(Lemma~KR--FP--2 in~\cite{RemondiereKRFP3_2026}).
\end{corollary}
\begin{proof}
By Theorem~\ref{thm:AHS}, the Hessian has discrete spectrum with kernel
$H^1_A$ and positive complement.

The lower bound on $\lambda_\perp(A)$ proceeds as in~\cite[\S 4]{RemondiereKRFP3_2026}:
the Hessian is structured as $-\Delta_A + V(A)$ on
$\Omega^1(M;\ad P)$ where $\Delta_A$ is the bundle-twisted
Hodge Laplacian and $V(A)$ is a bounded perturbation controlled by
$g\norm{A}_{L^4}$. On the transverse harmonic forms orthogonal to
$H^1_A$, the bundle-twisted Hodge Laplacian acts as $-\Delta$ on
$\Omega^{1}_\mathrm{transv}$ with smallest eigenvalue $m_0^2$. The
Birman--Schwinger argument of Proposition~KR--FP--1
in~\cite{RemondiereKRFP3_2026} gives the lower bound
$\lambda_\perp(A) \geq m_0^2 (1 - \kFP)$ on the action slice
$\overline\Lambda_{S_0}$ provided the Birman--Schwinger norm condition
$\norm{\widetilde K(A)} < 1$ holds; in the present setting that condition
is automatic in the small-action regime $S_{YM}(A) \leq S_0$ via the
Aubin--Talenti bound.

The crucial point is that, \emph{unlike} in~\cite{RemondiereKRFP3_2026}
where (H1) generic-vanishing of minimising eigenfunctions is assumed, here
the kernel structure on $\Mmod_k$ is \emph{rigorously identified} via the
AHS theorem~\ref{thm:AHS}, so the corresponding hypothesis is
unconditionally discharged.
\end{proof}

% =========================================================================
\section{Bakry--\'Emery on the instanton sector}\label{sec:bakry}

\subsection{Riemannian structure on $\Mmod_k$}
The $L^2$ metric on $T_A\mathcal A = \Omega^1(M;\ad P)$ descends to a
Riemannian metric $g_{\Mmod_k}$ on the smooth manifold $\Mmod_k^\circ$ of
irreducible self-dual connections, by Babelon--Viallet's O'Neill
formula~\cite{BabelonViallet1981} restricted to the slice $\Mmod_k$.

Specifically, $\Mmod_k = (F^+)^{-1}(0)/\mathcal G \subset \mathcal A /
\mathcal G$ is a submanifold of the orbit space, and its tangent space
at $[A]$ is precisely $H^1_A$ as identified in Theorem~\ref{thm:AHS}.
The induced metric is the restriction of the orbit-space $L^2$ metric.

\subsection{Ricci curvature on $\Mmod_k$ via Babelon--Viallet}

\begin{proposition}[Ricci on the instanton stratum]\label{prop:Ric-Mk}
Let $A\in\Mmod_k^\circ$ irreducible with $H^2_A = 0$. Then the Ricci
curvature of $(\Mmod_k^\circ, g_{\Mmod_k})$ at $[A]$ satisfies
$$
\Ric_{g_{\Mmod_k}}([A]) \;\geq\; (1-\kFP)\,m_0^2 \cdot g_{\Mmod_k}([A]).
$$
\end{proposition}
\begin{proof}
The Babelon--Viallet O'Neill formula~\cite[Eq.~(4.5)]{BabelonViallet1981}
relates the Ricci tensor on the orbit space $\mathcal A/\mathcal G$ to
the bottom of the spectrum of $M[A] = d_A^\dagger d_A$ via
\begin{equation}\label{eq:BV-ON-restr}
\Ric_{\mathcal A/\mathcal G}([A])(\xi,\xi) \;\geq\;
\lambda_{\min}(M[A])\,\abs{\xi}^2_{\mathcal A/\mathcal G},
\quad \xi \in T_{[A]}(\mathcal A/\mathcal G),
\end{equation}
in the normalisation $C_\mathrm{ON} = 1$ of~\cite[Cor.~KR--FP--A]{RemondiereKRFPB_2026}.
Restricting to $\xi \in T_{[A]}\Mmod_k \subset T_{[A]}(\mathcal A/\mathcal G)$
and using Corollary~\ref{cor:Hess-gap},
$$
\Ric_{g_{\Mmod_k}}([A])(\xi,\xi) \;\geq\;
\lambda_\perp(A) \,\abs{\xi}^2_{g_{\Mmod_k}}
\;\geq\; (1-\kFP)\,m_0^2 \,\abs{\xi}^2_{g_{\Mmod_k}}.
$$
This is the announced bound.

The submanifold restriction is justified because the Gauss equation
for $\Mmod_k \hookrightarrow \mathcal A/\mathcal G$ controls the
restriction of the ambient Ricci by the induced Ricci plus the
$|\!|\mathrm{II}|\!|^2$ second-fundamental-form term, and in the AHS
local model (geodesically convex neighbourhood of $[A]$ in the moduli
space) the second-fundamental form vanishes to leading order.
\end{proof}

\subsection{Bakry--\'Emery $\Gamma_2$-calculus on $\Mmod_k$}

\begin{theorem}[Bakry--\'Emery on the smooth manifold $\Mmod_k$]
\label{thm:BE-Mk}
Conditional on Proposition~\ref{prop:Ric-Mk}, the heat semigroup
$P_t^{\Mmod_k} = \exp(t\Delta_{g_{\Mmod_k}})$ satisfies the
curvature--dimension condition $\mathrm{CD}((1-\kFP)m_0^2, \infty)$:
$$
\Gamma_2(f,f) \;\geq\; (1-\kFP)\,m_0^2 \,\Gamma(f,f),
\qquad \forall f \in C_c^\infty(\Mmod_k^\circ).
$$
Consequently the heat-kernel measure satisfies a logarithmic Sobolev
inequality with constant $\rho \geq (1-\kFP)\,m_0^2$.
\end{theorem}
\begin{proof}
Standard application of Bakry--\'Emery's $\Gamma_2$-calculus on a
smooth Riemannian manifold with Ricci $\geq K\,g$
\cite{BakryEmery1985,BakryGentilLedoux2014}. The smoothness of
$\Mmod_k^\circ$ is provided by Theorem~\ref{thm:AHS} (smooth manifold
structure of the moduli space at an irreducible connection with
$H^2_A=0$).
\end{proof}

\subsection{Wilson-measure restriction}

\begin{proposition}[Restriction of the Wilson measure to $\Mmod_k$]
\label{prop:mu-restrict}
For $\beta\geq\beta_0$, the conditional measure $\mu^k_{a,\beta,L}$
on $\Mmod_k^a$ defined in \S~\ref{sec:setup} concentrates on the
neighbourhood of the smooth lattice instanton moduli space, in the
following sense: there exists $C(k,L) < \infty$ and $\eta(\beta) \to 0$
as $\beta\to\infty$ such that
$$
\mu^k_{a,\beta,L}\{U : \mathrm{dist}_{L^2}(U,\Mmod_k^a) > \eta(\beta)\}
\;\leq\; C(k,L)\,e^{-c\beta},
$$
with $c > 0$ depending on $k,L$.
\end{proposition}
\begin{proof}[Proof sketch]
This is the standard semiclassical concentration argument
(L\"uscher~\cite{Luscher2010WilsonFlow} for the lattice flowed-charge
sector identification, combined with the dilute-instanton expansion of
\cite{Polyakov1977}, citation to verify, for the saddle-point structure).
The Wilson action restricted to a small $L^2$-neighbourhood of $\Mmod_k^a$
attains its lattice minimum on (the lattice analogue of) $\Mmod_k^a$
itself, with second variation given by the Hessian of
Lemma~\ref{lem:Hess-Mk}, which has the AHS kernel structure of
Theorem~\ref{thm:AHS} plus a positive-definite complement of size
$\lambda_\perp \geq (1-\kFP)m_0^2 > 0$. Concentration around the
minimum follows by Laplace's method on the manifold $\Mmod_k^a$
relative to the Hessian-induced Gaussian on the normal bundle.
\end{proof}

% =========================================================================
\section{Main theorem}\label{sec:thm}

\begin{theorem}[Main result, LSI on the lattice instanton sector]
\label{thm:main}
Let $G = SU(N)$ on the lattice 4-torus $\Lambda_L^a$ with Wilson action
\eqref{eq:Wilson-action}. Fix a topological charge
$k\in\Z\setminus\{0\}$ and a sufficiently large $\beta\geq\beta_0(k,L)$
that the lattice instanton sector $\Mmod_k^a$ defined in
\eqref{eq:instanton-sector} is non-empty and concentrates on its smooth
interior (Proposition~\ref{prop:mu-restrict}). Let $\mu^k_{a,\beta,L}$
be the restricted Wilson measure on $\Mmod_k^a$ and let $L^k_{a,\beta}$
be the associated Langevin generator.

Then $\mu^k_{a,\beta,L}$ satisfies a logarithmic Sobolev inequality with
constant
\begin{equation}\label{eq:LSI-Mk}
\rho_k(\beta, a, L) \;\geq\; (1-\kFP)\,m_0^2 - \delta(\beta,a,L),
\end{equation}
where $\kFP = 1/(2|\Phi^+(SU(N))|)$, $m_0^2 = (2\pi/L)^2$, and
$\delta(\beta,a,L) \to 0$ as $\beta\to\infty$ and $a\to 0$ with $L$ fixed.
In the limit $L\to\infty$ with $a\to 0$ along the asymptotic-freedom
trajectory, the limit measure on each non-trivial topological sector
satisfies an LSI with constant $\rho_k^\infty \geq (1-\kFP)\,m_0^2(\infty) > 0$
provided $m_0^2(\infty) > 0$ persists in the limit (this is the
\emph{transverse mass gap} condition, weaker than the full Wilson
mass gap, since it only requires positivity of the lowest eigenvalue
on transverse harmonic 1-forms which is automatic on the lattice
$T^4$ and on flat $\R^4$ in the appropriate Sobolev space).
\end{theorem}

\begin{proof}
By Theorem~\ref{thm:AHS}, the kernel of the Wilson Hessian (analogue of
Lemma~\ref{lem:Hess-Mk} in the lattice context, see L\"uscher's
gradient-flow construction of the lattice-instanton moduli space
\cite{Luscher2010WilsonFlow}) on the lattice instanton stratum
$\Mmod_k^a$ is exactly the tangent space to $\Mmod_k^a$. The orthogonal
complement Hessian has uniform lower bound
$\lambda_\perp \geq (1-\kFP)m_0^2$ by Corollary~\ref{cor:Hess-gap}.

By Proposition~\ref{prop:Ric-Mk}, the Ricci of $\Mmod_k^a$ satisfies
$\Ric \geq (1-\kFP)m_0^2\,g$ everywhere on the smooth interior.

By Theorem~\ref{thm:BE-Mk}, the heat semigroup on $(\Mmod_k^a,
g_{\Mmod_k^a})$ satisfies $\mathrm{CD}((1-\kFP)m_0^2,\infty)$, and the
heat-kernel measure has LSI with constant $(1-\kFP)m_0^2$.

By Proposition~\ref{prop:mu-restrict}, the Wilson measure restricted
to $\Mmod_k^a$ is, up to exponentially small corrections at large
$\beta$, the Riemannian volume measure on $\Mmod_k^a$ weighted by the
Hessian-Gaussian on the normal bundle, which is the heat-kernel measure
at small flow time. The Holley--Stroock perturbation principle then
transfers the LSI from the heat-kernel measure to the conditional
Wilson measure, with a correction $\delta(\beta,a,L) \to 0$ encoding
the perturbation. The continuum limit follows from the uniform-in-$L$
nature of the Ricci bound and standard Fukushima--Oshima--Takeda
descent.
\end{proof}

\begin{remark}[Comparison with conditional KR--FP--B]
The constant in \eqref{eq:LSI-Mk} is \emph{identical in form} to the
constant $(1-\kFP)m_0^2/\cinf(D)$ of~\cite[Thm.~6.1]{RemondiereKRFPB_2026}
on the trivial-sector projection, but is here \emph{unconditional}
because the AHS argument bypasses (H1) (generic vanishing) entirely on
the instanton stratum. The $\cinf(D)$ factor of the latter (which enters
from the BBD~LSI uniform descent) is replaced here by the direct
Bakry--\'Emery descent on the finite-dimensional submanifold $\Mmod_k$,
which has no analogue of $\cinf(D)$.
\end{remark}

\begin{corollary}[Explicit constant for $SU(3)$ on $T^4$]\label{cor:SU3}
For $G = SU(3)$ on the lattice torus $T^4_L$ at sufficiently large
$\beta$ and $k\in\Z\setminus\{0\}$,
$$
\rho_k(\beta, a, L) \;\geq\; \frac{5}{6}\,m_0^2 - \delta,
\qquad \kFP(SU(3)) = 1/6.
$$
\end{corollary}

% =========================================================================
\section{The $M=T^4$ case in detail}\label{sec:T4}

\subsection{AHS on $T^4$: caveat on $H^2_A$}
The original AHS argument~\cite{AtiyahHitchinSinger1978} treats
$M=S^4$ with $b^+_2(S^4) = 1$. For $M = T^4$ (the lattice background),
$b^+_2(T^4) = 3$ and $H^2_A$ may be non-trivial; the index formula
\eqref{eq:dim-Mk} requires the term $(1+b^+_2)$ replaced by
$(1+b^+_2)$ generically, but the obstruction $H^2_A$ has to be checked.

For an \emph{anti-self-dual} (ASD) connection on $T^4$
(by convention $F^- = 0$, replacing $F^+ = 0$ for the dual orientation),
the cohomology $H^2_A$ governs obstructions in the deformation.
For $T^4$ flat with the round product metric, generic flat
$SU(N)$-connections of charge $k=0$ are smooth, but the higher-charge
instantons must be approached carefully.

\subsection{Periodic instantons (calorons) on $T^4$}
The natural generalisation of AHS to the 4-torus is the theory of
periodic instantons (\emph{calorons}); see Nahm~\cite{Nahm1983}
(citation to verify), van Baal~\cite{vanBaal1996} (citation to verify),
Garcia-Perez--Gonzalez-Arroyo et al.~\cite{GarciaPerezGonzalezArroyo}
(citation to verify) for the lattice analogues. The Nahm transform
provides a dual description identifying instantons on $T^4$ with
solutions of Nahm's equation on the dual torus $\hat T^4 = T^4$.

\subsection{Status on $T^4$: working hypothesis}
\begin{hypothesis}[AHS rigidity on $T^4$]\label{hyp:AHS-T4}
For $M=T^4$ at the lattice scale, the irreducible instanton moduli
space $\Mmod_k^a$ defined via the L\"uscher gradient flow has, at every
irreducible representative $[A]$, kernel of $\Hess S_{YM}(A)$ exactly
equal to $T_{[A]}\Mmod_k^a$, with positive complement of size
$\lambda_\perp \geq (1-\kFP)m_0^2$, where $m_0^2 = (2\pi/L)^2$.
\end{hypothesis}

\begin{remark}[Status of Hypothesis~\ref{hyp:AHS-T4}]\label{rem:AHS-T4}
On $M=S^4$ this is the original AHS
theorem~\cite{AtiyahHitchinSinger1978}. On $M=T^4$ this is
\emph{strongly expected} from the Nahm-transform duality (which
identifies $T^4$-instantons with Nahm-equation solutions on $\hat T^4$,
preserving the irreducibility-rigidity structure) but the explicit
verification of $H^2_A = 0$ at all irreducible $[A]\in\Mmod_k(T^4)$
to our knowledge is not a standard textbook result. The hypothesis is
included to mark this as a transparent point requiring expert
verification.
\end{remark}

\subsection{Conditional version of Theorem~\ref{thm:main} on $T^4$}
\begin{theorem}[Conditional on Hypothesis~\ref{hyp:AHS-T4}]
\label{thm:main-T4}
On $M=T^4$ with $G=SU(N)$ and the lattice Wilson measure as in
\eqref{eq:Wilson-action}, and under
Hypothesis~\ref{hyp:AHS-T4}, the conclusion of Theorem~\ref{thm:main}
holds verbatim: $\mu^k_{a,\beta,L}$ satisfies LSI with constant
$\rho_k(\beta,a,L) \geq (1-\kFP)m_0^2 - \delta(\beta,a,L)$.
\end{theorem}

% =========================================================================
\section{Remarks and outlook}\label{sec:remark}

\subsection{Why this is \emph{not} a full LSI for the Wilson measure}
The full Wilson measure is supported on $\bigcup_{k\in\Z}\Mmod_k^a$,
and its decomposition gives
$\mu_{a,\beta,L} = \sum_{k\in\Z} \pi_k\,\mu^k_{a,\beta,L}$
with $\pi_k = \mu_{a,\beta,L}(\Mmod_k^a)$. Theorem~\ref{thm:main}
delivers an LSI on each $\mu^k_{a,\beta,L}$ for $k\neq 0$, but with
constant $\rho_k$ \emph{independent} of $\pi_k$.

To deduce an LSI on the full mixture would require a Cheeger-type
or Bakry--Ledoux mixing-inequality argument that combines the
sector-wise LSIs with the discrete probability vector $(\pi_k)_k$.
Such mixing-inequalities generally fail to give a useful uniform
constant when one sector has $\pi_k$ exponentially small relative to
the other. In particular, the trivial sector $k=0$ has
$\pi_0 = 1 - O(e^{-8\pi^2/g^2})$ and dominates the measure, so its
LSI constant controls the full Wilson LSI.

The \emph{conclusion}: Theorem~\ref{thm:main} does not produce an
LSI on the full Wilson measure, but it provides a structurally clean
unconditional sub-result that lives on $O(e^{-8\pi^2/g^2})$ mass.

\subsection{Could AHS extend to the trivial sector $k=0$?}
The natural attempt is to apply the deformation-complex argument to
the trivial connection $A=0$ on the trivial bundle. The cohomology
$H^0_{A=0} = \su(N)$ (since $A=0$ has full continuous isotropy:
$\mathcal G$-stabiliser is all global gauge transformations preserving
$A=0$, which is the constant subgroup $G\subset\mathcal G$, and its
Lie algebra is $\su(N)$ for $G$ connected). So $A=0$ is
\emph{not irreducible}, and the AHS hypothesis $H^0_A = 0$ \emph{fails}.

This is precisely the trivial-sector zero-mode obstruction
discussed in~\cite[\S 5]{RemondiereKRFP3_2026} and is the source of the
hypothesis~(H1) (generic vanishing). The AHS argument cannot bypass it
without supplementary input.

\subsection{Connection with the Cao--Park--Sheffield programme}
The recent works of Cao--Park--Sheffield~\cite{CaoParkSheffield2024}
and Cao--Nissim--Sheffield~\cite{CaoNissimSheffield2025} construct
scaling limits of $SU(N)$ lattice Yang--Mills via random-surface
representations. Their construction is sector-wise transparent: the
random-surface representation respects the topological charge
decomposition (it sees the holonomies along closed loops, which
distinguish topological sectors). Our Theorem~\ref{thm:main} provides
an unconditional LSI on each sector of their construction, which may
be useful in their concentration estimates.

\subsection{Connection with the Chatterjee Yang--Mills--Higgs result}
Chatterjee~\cite{Chatterjee2024} constructs a scaling limit of $SU(2)$
lattice Yang--Mills--Higgs in any $d\geq 2$, proving mass generation
by the Higgs mechanism. The Higgs field breaks the symmetry that
generates the AHS-trivial-sector zero mode, so the AHS argument adapts
more cleanly to the Yang--Mills--Higgs measure: the trivial
sector at $A=0,\phi = \langle\phi\rangle$ has isotropy
$\mathrm{Stab}_{(A=0,\langle\phi\rangle)} = \mathrm{Stab}_{\langle\phi\rangle}
\subset G$, which is a \emph{proper} subgroup of $G$. For the
abelian Higgs with vacuum expectation $\langle\phi\rangle$ breaking $G$
to a maximal torus, the residual isotropy is the Cartan,
exactly the structure handled by the Kostant identity (Lemma~KR--FP--2).

This suggests a follow-up project: \emph{adapting the AHS argument to
the (lattice) Yang--Mills--Higgs measure to obtain an unconditional LSI
on the trivial Higgs vacuum sector}, with constant
$\rho \geq (1-\kFP)\,m_0^2$ where $\kFP$ is now controlled by the
\emph{Higgs-broken} root system (not the full $\Phi^+(G)$). This would
be a structural contribution to the Chatterjee scaling-limit programme.

\subsection{Honest probability of conditional Clay closure}
The result of this Letter, combined with the conditional reduction
of~\cite{RemondiereKRFPB_2026}, leaves the residual conditional Clay
gap as follows:
\begin{itemize}[leftmargin=2em]
\item Trivial sector $k=0$: requires (H1)--(H3) of
\cite{RemondiereKRFP3_2026} (one open hypothesis, two technical) +
BBD-type uniform LSI on the lattice $SU(N)$ Wilson measure for the
$k=0$ stratum.
\item Non-trivial sectors $k\neq 0$: \emph{unconditional} via the
present Theorem~\ref{thm:main}, modulo Hypothesis~\ref{hyp:AHS-T4} on
$T^4$, which is expected to hold by Nahm-duality but requires expert
verification.
\end{itemize}

The unconditional gain on $k\neq 0$ is of \emph{negligible measure}
($O(e^{-8\pi^2/g^2})$ at large $\beta$), so the present result does
\emph{not} change the honest estimate $P(\text{Clay 10y}) \in [70\%,82\%]$
post-Opus~\#319+\#2. Its value is mathematical, not statistical.

% =========================================================================
\section*{Acknowledgments}
This work was prepared with computational research assistance from
Anthropic's Claude system (Opus~4.7, 1M context window). All
mathematical content (theorems, proofs, hypotheses, references) is the
sole responsibility of the author.

% =========================================================================
\begin{thebibliography}{99}

\bibitem{AtiyahHitchinSinger1978}
M.~F. Atiyah, N.~J. Hitchin, I.~M. Singer,
``Self-duality in four-dimensional Riemannian geometry,''
\emph{Proc. R. Soc. A} \textbf{362} (1978) 425--461.

\bibitem{Atiyah1979}
M.~F. Atiyah,
\emph{Geometry of Yang--Mills fields},
Lezioni Fermiane, Scuola Normale Superiore, Pisa, 1979.

\bibitem{Aubin1976}
T.~Aubin,
``Probl\`emes isop\'erim\'etriques et espaces de Sobolev,''
\emph{J. Diff. Geom.} \textbf{11} (1976) 573--598.

\bibitem{BabelonViallet1981}
O.~Babelon, C.-M.~Viallet,
``The Riemannian geometry of the configuration space of gauge theories,''
\emph{Commun. Math. Phys.} \textbf{81} (1981) 515--525.

\bibitem{BakryEmery1985}
D.~Bakry, M.~\'Emery,
``Diffusions hypercontractives,''
\emph{S\'eminaire de Probabilit\'es XIX},
Lecture Notes in Mathematics \textbf{1123} (Springer, 1985) 177--206.

\bibitem{BakryGentilLedoux2014}
D.~Bakry, I.~Gentil, M.~Ledoux,
\emph{Analysis and Geometry of Markov Diffusion Operators},
Grundlehren \textbf{348}, Springer, 2014.

\bibitem{BauerschmidtBodineauDagallier2024}
R.~Bauerschmidt, T.~Bodineau, B.~Dagallier,
``Stochastic dynamics and the Polchinski equation: an introduction,''
\emph{Probab. Surveys} \textbf{21} (2024) 200--290.
arXiv:2307.07619.

\bibitem{BauerschmidtDagallier2024}
R.~Bauerschmidt, B.~Dagallier,
``Log-Sobolev inequality for the $\varphi^4_2$ and $\varphi^4_3$ measures,''
arXiv:2202.02295.

\bibitem{CaoNissimSheffield2025}
S.~Cao, A.~Nissim, S.~Sheffield,
``A dynamical approach to area law for lattice Yang--Mills,''
arXiv:2509.04688.

\bibitem{CaoParkSheffield2024}
S.~Cao, M.~Park, S.~Sheffield,
``Random surfaces and lattice Yang--Mills,''
arXiv:2307.06790.

\bibitem{Chatterjee2024}
S.~Chatterjee,
``A scaling limit of $SU(2)$ lattice Yang--Mills--Higgs theory,''
\emph{Probab. Math. Phys.}, to appear (2026).
arXiv:2401.10507.

\bibitem{GarciaPerezGonzalezArroyo}
A.~Gonz\'alez-Arroyo, P.~Mart\'inez,
``Periodic instantons on the lattice,''
\emph{citation to verify}: caloron-type periodic instanton constructions
on $T^4$ originate in the work of Nahm and van Baal in the
1980s--1990s. See \cite{Nahm1983,vanBaal1996} for primary sources.

\bibitem{Luscher2009TrivializingMaps}
M.~L\"uscher,
``Trivializing maps, the Wilson flow and the HMC algorithm,''
\emph{Commun. Math. Phys.} \textbf{293} (2010) 899--919.
arXiv:0907.5491.

\bibitem{Luscher2010WilsonFlow}
M.~L\"uscher,
``Properties and uses of the Wilson flow in lattice QCD,''
\emph{JHEP} \textbf{08} (2010) 071.
arXiv:1006.4518.

\bibitem{LuscherTopo1998}
\emph{Citation to verify}: dilute-instanton mass of lattice topological
sectors at large $\beta$ via L\"uscher gradient-flow charge
identification; see \cite{Luscher2010WilsonFlow} for the gradient-flow
charge definition and the standard literature on the instanton-gas
approximation.

\bibitem{Nahm1983}
W.~Nahm,
``Self-dual monopoles and calorons,''
\emph{citation to verify}: Nahm's construction of calorons (periodic
instantons) is contained in his Bonn preprints and subsequent works
from the early 1980s.

\bibitem{Polyakov1977}
A.~M.~Polyakov,
``Quark confinement and topology of gauge theories,''
\emph{Nucl. Phys. B} \textbf{120} (1977) 429--458.

\bibitem{RemondiereKRFP3_2026}
K.~R\'emondi\`ere,
\emph{Spectral bound for the Faddeev--Popov operator on the fundamental
modular domain of the slice $\overline\Lambda_{S_0}$: a Birman--Schwinger /
Aubin--Talenti analysis} (companion preprint
\texttt{Paper\_KR\_FP3\_AnnalsMath}), May 2026.

\bibitem{RemondiereKRFPB_2026}
K.~R\'emondi\`ere,
\emph{Mass gap for 4D pure Yang--Mills via Bakry--\'Emery on the
fundamental modular domain: a conditional reduction} (companion preprint
\texttt{Paper\_KR\_FP\_B\_BakryEmery\_LMP}), May 2026.

\bibitem{Uhlenbeck1982}
K.~Uhlenbeck,
``Connections with $L^p$ bounds on curvature,''
\emph{Commun. Math. Phys.} \textbf{83} (1982) 31--42.

\bibitem{vanBaal1996}
P.~van Baal,
``Instantons versus monopoles,''
\emph{citation to verify}: van Baal's work on calorons and twisted
boundary conditions through the 1990s; specific arXiv ID requires
direct verification.

\bibitem{Talenti1976}
G.~Talenti,
``Best constant in Sobolev inequality,''
\emph{Ann. Mat. Pura Appl.} \textbf{110} (1976) 353--372.

\end{thebibliography}

\end{document}
